\section{Primary birth regions and adjacent fine meshes}
\label{sec:area_law_primary_regions}

Retain the half-open cells $C_{k,z}$, the layer $D_k$, and its finite
index set $I_k$ from Definition~\ref{def:area_law_dyadic_layers}.
The origin $o\in\mathbb R^2$ is arbitrary, the endpoint set
$Z\subset\mathbb Z^2$ is finite, and the neighborhood radius $C_0$ is
a nonnegative integer. All distances and diameters use the sup metric
on $\mathbb R^2$.

The primary regions in \cite[Section~11]{OpenAI2026AreaLaw} are obtained
by intersecting a pitch interior with an actual layer and taking its
closure. All fragments in one pitch interior receive the same primary
identifier, even when their union is disconnected.

% Source: 10-geometry.tex, lines 212--218 and 237--249,
% geometry:primary-pieces. These entries describe individual primary
% birth regions and their fragments; they do not state the full partition
% or the scanner-template conditions of prop:two-families.

\begin{definition}[Pitch interiors and primary birth regions]
    \label{def:area_law_primary_regions}
    \lean{TNLean.PEPS.AreaLaw.Geometry.pitchInterior,
        TNLean.PEPS.AreaLaw.Geometry.primaryFragment,
        TNLean.PEPS.AreaLaw.Geometry.primaryFragmentIndices,
        TNLean.PEPS.AreaLaw.Geometry.primaryBirthRegion}
    \leanok
    \uses{def:area_law_dyadic_cells,def:area_law_dyadic_layers}
    Fix nonnegative integer indices $k,\ell,p$, put
    $r=2^k$, $t=2^\ell$, and $s=2^p$, and choose integer shifts $a,b$.
    For every $j\in\mathbb Z^2$, set
    \begin{align}
        v_j     & =o+t(a,b)+sj,                                \\
        U_j     & =(v_{j,1}+t,v_{j,1}+s)
        \times(v_{j,2}+t,v_{j,2}+s),                           \\
        F_{j,z} & =\overline{C_{k,z}\cap U_j},                 \\
        J_j     & =\{z\in I_k:C_{k,z}\cap U_j\ne\varnothing\}, \\
        R_j     & =\overline{D_k\cap U_j}.
    \end{align}
    Here $U_j$ is the open pitch interior, $F_{j,z}$ is a closed cell
    fragment, and $R_j$ is the corresponding primary birth region.
    These definitions also allow empty interiors and empty fragments.
    In the geometric construction, $\ell=\ell_k$, $p=p_k$, and
    $a,b$ represent the selected residue classes modulo $s/t$.
\end{definition}

\begin{theorem}[Finite rectangular decomposition]
    \label{thm:area_law_primary_fragment_union}
    \lean{TNLean.PEPS.AreaLaw.Geometry.primaryBirthRegion_eq_iUnion_primaryFragment}
    \leanok
    \uses{def:area_law_primary_regions}
    For every choice of the preceding data,
    \begin{align}
        R_j & =\bigcup_{z\in J_j}F_{j,z}.
    \end{align}
\end{theorem}
\begin{proof}\leanok
    \uses{thm:area_law_dyadic_nesting}
    The exact cell union gives $D_k=\bigcup_{z\in I_k}C_{k,z}$.
    Intersect this finite union with $U_j$ and commute closure with
    the finite union. Empty intersections have empty closures, so
    retaining precisely the indices in $J_j$ leaves the union unchanged.
\end{proof}

\begin{theorem}[Closed cell fragments]
    \label{thm:area_law_primary_fragment_rectangles}
    \lean{TNLean.PEPS.AreaLaw.Geometry.primaryFragment_eq_closedRectangle,
        TNLean.PEPS.AreaLaw.Geometry.primaryFragment_subset_closure_dyadicCell,
        TNLean.PEPS.AreaLaw.Geometry.primaryFragment_dist_le,
        TNLean.PEPS.AreaLaw.Geometry.primaryFragment_diam_le}
    \leanok
    \uses{def:area_law_primary_regions}
    For $i\in\{1,2\}$, write
    $\alpha_{i,z}=o_i+rz_i$ and $\beta_{i,z}=o_i+r(z_i+1)$.
    If $C_{k,z}\cap U_j\ne\varnothing$, then
    \begin{align}
        F_{j,z}
         & =\prod_{i=1}^2
        [\max\{\alpha_{i,z},v_{j,i}+t\},
        \min\{\beta_{i,z},v_{j,i}+s\}].
    \end{align}
    Both coordinate intervals have strictly positive length. Every
    fragment, including an empty one, satisfies the following containment
    and diameter bounds. The distance bound holds for arbitrary points
    $x,y\in F_{j,z}$.
    \begin{align}
        F_{j,z}                            & \subseteq\overline{C_{k,z}}, \\
        \|x-y\|_\infty                     & \le r,                       \\
        \operatorname{diam}_\infty F_{j,z} & \le r.
    \end{align}
\end{theorem}
\begin{proof}\leanok
    \uses{thm:area_law_dyadic_cell_distance}
    A point of the nonempty intersection lies below both upper endpoints
    in each coordinate. It lies at or above the cell's lower endpoint
    and strictly above the pitch interior's lower endpoint. Hence the
    maximum lower endpoint is strictly below the minimum upper endpoint.
    The coordinate intersection contains the open interval between these
    endpoints and is contained in their closed interval. Taking closures
    therefore gives that closed interval. Closure commutes with products,
    proving the rectangle formula.

    Containment in the closed cell follows from monotonicity of closure.
    Apply the closed-cell distance bound with identical cell indices and
    index radius zero. This gives the pointwise distance bound, and hence
    the diameter bound.
\end{proof}

\begin{theorem}[Primary diameter and separation]
    \label{thm:area_law_primary_diameter_separation}
    \lean{TNLean.PEPS.AreaLaw.Geometry.primaryBirthRegion_subset_closure_pitchInterior,
        TNLean.PEPS.AreaLaw.Geometry.primaryBirthRegion_dist_le,
        TNLean.PEPS.AreaLaw.Geometry.primaryBirthRegion_diam_le,
        TNLean.PEPS.AreaLaw.Geometry.primaryBirthRegion_dist_separation}
    \leanok
    \uses{def:area_law_primary_regions}
    For every primary index $j$, the following containment and diameter
    bounds hold. The distance bound holds for arbitrary points $x,y\in R_j$.
    \begin{align}
        R_j                            & \subseteq\overline{U_j}, \\
        \|x-y\|_\infty                 & \le s,                   \\
        \operatorname{diam}_\infty R_j & \le s.
    \end{align}
    For distinct indices $i,j$, every $x\in R_i$ and $y\in R_j$ satisfy
    \begin{align}
        \|x-y\|_\infty & \ge t.
    \end{align}
\end{theorem}
\begin{proof}\leanok
    \uses{def:area_law_primary_regions}
    Monotonicity of closure gives $R_j\subseteq\overline{U_j}$.
    This closure is contained in the product of the closed intervals
    $[v_{j,q}+t,v_{j,q}+s]$, for $q=1,2$. Two points in this product
    have coordinate differences at most $s$, proving the distance and
    diameter bounds.

    Distinct integer indices differ in at least one coordinate. Suppose
    $i_q<j_q$. The upper endpoint for index $i_q$ is at most
    $o_q+t(a,b)_q+sj_q$, while the lower endpoint for index $j_q$ is
    $o_q+t(a,b)_q+sj_q+t$. Thus the two coordinate intervals are
    separated by at least $t$. Taking the maximum coordinate distance
    proves the separation claim; the other ordering is symmetric.
\end{proof}

\begin{theorem}[Number of fragments of one primary]
    \label{thm:area_law_primary_fragment_count}
    \lean{TNLean.PEPS.AreaLaw.Geometry.card_primaryFragmentIndices_le}
    \leanok
    \uses{def:area_law_primary_regions}
    If $k\le p$, then every primary index satisfies
    \begin{align}
        |J_j| & \le\bigl(2^{p-k}+2\bigr)^2
        =\bigl(s/r+2\bigr)^2.
    \end{align}
    This bound is independent of the endpoint set, neighborhood radius,
    origin, and selected shifts.
\end{theorem}
\begin{proof}\leanok
    \uses{def:area_law_primary_regions}
    Put $N=2^{p-k}$, so $s=Nr$, and let
    $q_i=\lfloor(v_{j,i}-o_i)/r\rfloor$. If a cell indexed by $z$ meets
    $U_j$, a point in that intersection lies above $v_{j,i}$ and below
    $v_{j,i}+s$. The half-open cell inequalities and the defining floor
    inequalities consequently give, for each coordinate $i\in\{1,2\}$,
    \begin{align}
        q_i & \le z_i\le q_i+N.
    \end{align}
    Thus $J_j$ lies in an integer rectangle with $N+1$ indices in each
    coordinate. Its cardinality is at most $(N+1)^2$, which is at most
    $(N+2)^2$.
\end{proof}

% Source: 10-geometry.tex, lines 299--306 and 352--356,
% the adjacent-mesh arithmetic used in geometry:initial-stars.
\begin{theorem}[Adjacent fine meshes]
    \label{thm:area_law_adjacent_fine_meshes}
    \lean{TNLean.PEPS.AreaLaw.Geometry.fineScaleIndex_mono,
        TNLean.PEPS.AreaLaw.Geometry.fineScaleIndex_succ,
        TNLean.PEPS.AreaLaw.Geometry.fineScale_side_succ}
    \leanok
    \uses{def:area_law_fine_belt_scales}
    The fixed fine exponents $\ell_k=\lfloor\zeta k\rfloor$ are
    nondecreasing. For every nonnegative integer $k$,
    \begin{align}
        \ell_{k+1} & \in\{\ell_k,\ell_k+1\}, \\
        t_{k+1}    & \in\{t_k,2t_k\}.
    \end{align}
\end{theorem}
\begin{proof}\leanok
    \uses{thm:area_law_exponent_gaps}
    Since $0\le\zeta\le1$, multiplication by $\zeta$ and taking floors
    preserve the order of the layer indices. Also
    $\zeta(k+1)\le\zeta k+1$, and adding one commutes with taking the
    floor. Therefore $\ell_k\le\ell_{k+1}\le\ell_k+1$.
    Exponentiating the two possible integer values gives the side-length
    assertion.
\end{proof}
